\section{Jet separation and a nef-and-big decomposition}
\label{jet:iitaka-section}

There is another way to obtain the boundary needed for nonvanishing.
A metric that is strictly positive on one open set first separates
linearly growing jets at a point.  This proves bigness directly by Nadel
vanishing.  A subsequent multiplier-ideal construction produces a nef
and big divisor on a resolution, which can be perturbed by any fixed
pseudoeffective class.  We prove this base lemma independently, then
apply it to an Iitaka fibration.  The Iitaka construction also provides
a different proof of the horizontal curvature bound.

\subsection{A base lemma proved by jets and a nef-and-big decomposition}

\begin{theorem}\label{jet:base-lemma}
Let $Y$ be a smooth connected projective variety of dimension $b>0$,
let $A\geq0$ be an integral divisor, and put $M=-K_Y+A$.
Assume $M$ has a singular Hermitian metric with psh local weights
$\varphi$, smooth and strictly positively curved on a nonempty open
subset, and satisfying
\begin{equation}\label{jet:weighted-assumption}
 e^{-\varphi}|s_A|^2\in L^{1+\gamma}_{\mathrm{loc}}
 \qquad\text{for some }\gamma>0.
\end{equation}
For every pseudoeffective $\mathbb Q$-divisor $W$ on $Y$ there is an
effective $\mathbb Q$-divisor $C_A$ supported on $\operatorname{Supp}A$
such that $W+C_A$ is $\mathbb Q$-linearly equivalent to an effective
$\mathbb Q$-divisor.
\end{theorem}

\begin{proof}
We first prove that $M$ is big, using its strict positivity only on the
specified patch.  Choose ample Cartier divisors $H_1,A_1$ so that
$B_1+A_1$ is very ample, where $B_1=K_Y+H_1$.
At a point $y_0$ in the patch, use coordinates $z$ and a smooth cutoff
to form a function $\ell$ equal to $\log\|z\|^2$ near $y_0$, supported
in a ball relatively compact in that patch, and zero outside it.
Fix a background K\"ahler form $\omega$.  On the support of the
cutoff error choose constants $\delta>0$ and $C>0$ such that
$i\partial\bar\partial\varphi\geq\delta\omega$ and
$i\partial\bar\partial\ell\geq-C\omega$; outside that support the
added logarithmic pole has nonnegative curvature.  Fix
$0<c\leq\delta/(2C)$.  For large integers $h$ and
$N=\lfloor ch\rfloor$, the metric
on $hM+H_1$ with weight
\[
 h\varphi+N\ell+\varphi_{H_1}
\]
therefore has curvature bounded below by a K\"ahler form.
Indeed the strict curvature of $h\varphi$ absorbs the cutoff error,
and the positive metric on $H_1$ supplies a global lower bound.

Let $\mathcal I_h$ be this metric's multiplier ideal.  Nadel vanishing
\cite{Nadel1990} gives
\[
 H^1\bigl(Y,\mathcal O_Y(hM+B_1)\otimes\mathcal I_h\bigr)=0.
\]
Near $y_0$, the original weight is smooth and
\[
 \mathcal I_{h,y_0}=\mathfrak m_{y_0}^{N-b+1}
\]
for $N\geq b$.  In fact a holomorphic monomial of total order $q$
is integrable against $\|z\|^{-2N}$ exactly when $q>N-b$.
The support of the quotient by this ideal has $y_0$ as an isolated
component, separated from its other support by a neighborhood on which
the original weight is smooth.  Any jet in
$\mathcal O_{Y,y_0}/\mathfrak m_{y_0}^{N-b+1}$ thus defines a global
section of this quotient, taken to be zero on the other components.
The vanishing lifts every such jet.  The quotient has length
$\binom{N}{b}$, so in fact
\[
 h^0(Y,hM+B_1)\geq\binom{\lfloor ch\rfloor}{b}.
\]
Consequently
\begin{equation}\label{jet:section-growth}
 h^0(Y,hM+B_1)\geq c' h^b
\end{equation}
for a constant $c'>0$ and all sufficiently large $h$.

Take a fixed smooth divisor $T_0\in|B_1+A_1|$.  We have
$h^0(T_0,(hM+B_1)|_{T_0})=O(h^{b-1})$.
For completeness, choose an ample divisor $H_2$ such that
$H_2-M$ has a section not identically zero on any component of $T_0$;
its $h$-th power embeds these restricted sections into those of
$(hH_2+B_1)|_{T_0}$.  The latter dimensions have the asserted growth by
the ample Hilbert polynomial, including the bounded zero-dimensional
case when $b=1$.  The restriction exact sequence and
\eqref{jet:section-growth} now imply
\[
 H^0(Y,hM-A_1)\ne0
\]
for large $h$.  Thus $M$ is big.

Write $M$ as the sum of an ample rational class and an effective
rational divisor.  This gives another metric with weight $\psi$,
curvature bounded below by a positive multiple of a K\"ahler form,
and divisorial analytic singularities.  Small inverse powers of its
local divisor equation are integrable, so
$e^{-\psi}|s_A|^2\in L^\tau_{\mathrm{loc}}$ for some $\tau>0$.
Replace $\varphi$ by $(1-\sigma)\varphi+\sigma\psi$, for a sufficiently
small $\sigma>0$.  Its curvature now has a global positive lower bound.
It also retains the local integrability
\begin{equation}\label{jet:mixed-l1}
 e^{-\varphi}|s_A|^2\in L^1_{\mathrm{loc}}.
\end{equation}
Indeed the new expression is the product of the two old weighted
expressions to the powers $1-\sigma$ and $\sigma$, and H\"older applies
as soon as
$(1-\sigma)/(1+\gamma)+\sigma/\tau<1$.
We keep the notation $\varphi$ for this new weight.

We now extract a nef-and-big divisor on a resolution, while retaining
the integrability needed to bound every discrepancy.
Fix a very ample divisor $H$ with a positive metric.  Take an integer
$k>0$ sufficiently large that $P=kM-H$ has a semipositive metric with
weight $k\varphi-\varphi_H$.  Put
\[
 B_0=K_Y+(b+1)H,
 \qquad C=mP+B_0,
\]
where $m$ is a positive integer so large that $mH-B_0$ is ample.
Set $\mathcal J=\mathcal I(mk\varphi)$.
For $1\leq j\leq b$, Nadel vanishing applied to
$mP+(b+1-j)H$ gives
\[
 H^j\bigl(Y,\mathcal O_Y(C-jH)\otimes\mathcal J\bigr)=0.
\]
The added smooth weights leave the multiplier ideal unchanged, and
the curvature is strictly positive because $b+1-j\geq1$.
Castelnuovo--Mumford regularity therefore makes
$\mathcal O_Y(C)\otimes\mathcal J$ globally generated.

Resolve $\mathcal J$ and $A$ by a smooth projective birational morphism
$d:Z\to Y$, including the exceptional divisor in the resolved normal
crossings support.  Write
$\mathcal J\mathcal O_Z=\mathcal O_Z(-G)$.
Then $d^*C-G$ is globally generated, and
\begin{equation}\label{jet:nef-big-decomposition}
 d^*M\sim_{\mathbb Q}Q+\frac{G}{mk},\qquad
 Q=\frac{d^*C-G+d^*(mH-B_0)}{mk}.
\end{equation}
The divisor $Q$ is nef and big: its first summand is nef, and the
pullback of the ample divisor $mH-B_0$ is nef and big.
This decomposition will allow us to perturb by $d^*W$ while preserving
control of the boundary coefficients.

Put
\begin{equation}\label{jet:delta-zero}
 \Delta_0=\frac{G}{mk}-K_{Z/Y}-d^*A.
\end{equation}
We claim that $\Delta_0^+$ is a klt boundary.  At a general point of
a prime $R=(z_1=0)$ in the resolved support, let $r,a,c$ be its
coefficients in $G,K_{Z/Y},d^*A$, respectively.
The Ohsawa--Takegoshi point-extension theorem
\cite{OhsawaTakegoshi1987} gives
\begin{equation}\label{jet:point-estimate}
 e^{mk\varphi(d(z))}\leq C_1|z_1|^{2r}.
\end{equation}
Explicitly, point extension gives holomorphic functions with uniformly
bounded $mk\varphi$-weighted $L^2$ norm and
$|q_{y'}(y')|^2\geq c_1e^{mk\varphi(y')}$ for a fixed $c_1>0$.  Local upper bounds for the weight give
ordinary sup bounds.  Pullback makes every such function divisible by
$z_1^r$; uniform Cauchy bounds on the quotient and $y'=d(z)$ give
\eqref{jet:point-estimate}.  This argument uses the coherent multiplier
ideal itself, and requires no analytic-singularity hypothesis on
$\varphi$.

Changing variables in \eqref{jet:mixed-l1} includes the squared
Jacobian of $d$, of order $2a$.  Inequality
\eqref{jet:point-estimate} then forces integrability of
$|z_1|^{2(a+c-r/(mk))}$.  Hence
\begin{equation}\label{jet:strict-discrepancy}
 \frac r{mk}<1+a+c.
\end{equation}
All coefficients of $\Delta_0$ are consequently less than one.
Its support has simple normal crossings, so $(Z,\Delta_0^+)$ is klt.
Moreover \eqref{jet:nef-big-decomposition} and $M=-K_Y+A$ imply
\[
 K_Z+\Delta_0\sim_{\mathbb Q}-Q.
\]

It remains to incorporate $W$ without losing the strict klt inequalities.
Write $Q\sim_{\mathbb Q}A_0+G'$ with $A_0$ ample and $G'\geq0$.
For a small rational $\lambda\in(0,1)$,
\[
 Q-\lambda G'\sim_{\mathbb Q}(1-\lambda)Q+\lambda A_0
\]
is ample.  Check the pair on a common log resolution of $\Delta_0^+$ and
$G'$.  Its crepant boundary coefficients depend affinely on $\lambda$
and are strictly below one at $\lambda=0$.  Thus
$(Z,\Delta_0^++\lambda G')$ remains klt for sufficiently small
$\lambda$.  This resolution is used only to check discrepancies;
$Q$ and $\Delta_0$ continue to denote the divisors on $Z$.  By openness of the ample cone choose a further
small rational $\epsilon>0$ so that
$Q-\lambda G'+\epsilon d^*W$ is ample.
Take a general effective rational representative $A'$ of this ample
class.  Using a sufficiently high very ample multiple dilutes its
coefficients; Bertini on the same resolution ensures that adding it
preserves klt and gives it no component in common with $\Delta_0$.
Set
\[
 B'=\lambda G'+A',\qquad \Delta=\Delta_0+B'.
\]
Then $\Delta^+\leq\Delta_0^++\lambda G'+A'$ is klt.  Its positive
part is big because it contains the ample rational divisor $A'$.
Finally
\[
 K_Z+\Delta^+\sim_{\mathbb Q}\epsilon d^*W+\Delta^-
\]
is pseudoeffective.  BCHM big-boundary nonvanishing
\cite[Theorem~D]{BCHM2010} gives real-linear effectivity, and
Lemma~\ref{lem:rational-recovery} gives rational linear effectivity.
Pushing forward and dividing by $\epsilon$ proves the theorem with
\[
 C_A=\epsilon^{-1}d_*\Delta^-.
\]
This divisor is effective and supported on $A$: the only negative
terms introduced in \eqref{jet:delta-zero} were $K_{Z/Y}$, which is
exceptional, and $d^*A$.
\end{proof}

\subsection{The Iitaka base revisited}

We apply the jet construction to the Iitaka base already used in
Section~\ref{conv:iitaka}.  The algebraic return remains
Lemma~\ref{conv:return}; the new point is a fiber-intersection proof
that the original smooth metric is strictly positive in base directions.
We also record the stronger generic-fiber rank statement behind this
version of the argument.

Retain $X$ and $L=-K_X$ with its smooth semipositive metric from
Theorem~\ref{conv:fixed-error}.  Write the given effective integral
divisors as $N_m\sim mL-D$ for an unbounded set of positive integers
$m$, with $D$ pseudoeffective Cartier, and fix two exponents
$j<q$.  Use the Iitaka fibration of $N_q$ constructed in
Section~\ref{conv:iitaka}.  That construction settles the cases of
zero-dimensional and generically finite base.  In the remaining case
it gives
\[
 X\xleftarrow{\ \pi\ }S\xrightarrow{\ f\ }Y,
 \qquad 0<b=\dim Y<n=\dim X,
 \qquad u\pi^*N_q\sim f^*A_0+I,
\]
where $A_0$ is ample Cartier, $u$ is a positive integer, and $I\ge0$.  The varieties are
smooth projective, $f$ has connected fibers, and every source prime
mapping into base codimension at least two is $\pi$-exceptional.
Every ratio from a multiple of $N_q$ belongs to $\mathbb C(Y)$, which
is relatively algebraically closed in $\mathbb C(X)$.
These are precisely the model and field properties established there;
no flatness of the final smooth resolution is needed.

Put $E=K_{S/X}$ and $D_S=\pi^*N_q+E$.  Adding positive multiples of
$E$ does not change the sections of multiples of $\pi^*N_q$, because
$E$ is exceptional and $X$ is normal.  Thus $\kappa(D_S)=b$ and
$uD_S-f^*A_0$ has an effective representative.  In fact
\begin{equation}\label{jet:all-ranks}
 \operatorname{rank}f_*\mathcal O_S(\ell D_S)=1
 \qquad(\ell\in\mathbb Z_{>0}).
\end{equation}
For if this rank exceeded one, an ample base twist would yield two
sections of $\ell D_S+af^*A_0$ independent over $\mathbb C(Y)$.
Multiplication by the section of $a(uD_S-f^*A_0)$ places them in
$(\ell+au)D_S$.  Their ratio is transcendental over $\mathbb C(Y)$.
Other sections of a sufficiently divisible multiple recover $b$
algebraically independent base functions through $A_0$.  Products put
all these ratios into one multiple of $D_S$, contradicting
$\kappa(D_S)=b$.  The canonical sections show that the rank cannot be
zero.  At the generic point $\eta$ of $Y$, this gives
$h^0(S_\eta,E|_{S_\eta})=1$, and
$\pi^*N_q|_{S_\eta}$ has Iitaka dimension zero.

The interpolation and vertical-minimum construction of
Section~\ref{conv:iitaka}, with these indices, gives
\[
 R=\frac{\pi^*N_q-\pi^*N_j}{q-j}\sim_{\mathbb Q}\pi^*L,
 \qquad
 W=\sum_T\min_{V\mapsto T}
       \frac{\operatorname{ord}_V R}{\operatorname{ord}_V f^*T}\,T.
\]
Here $R$ has nonnegative horizontal part and $W$ is pseudoeffective.
The rank statement above also gives the interpolation directly:
for each given $m>q$, restrict the two effective divisors
$(m-j)N_q\sim(m-q)N_j+(q-j)N_m$ to the generic fiber.  Its
one-dimensional space of sections makes their horizontal divisors
equal.  The horizontal slope is therefore nonnegative as $m$ tends
to infinity.  The principal difference has vertical divisor, so its
restriction to the projective normal generic fiber is constant; hence
it is pulled back from $Y$.  Taking vertical minima in this equality
recovers the pseudoeffectivity of $W$ by the fixed-error limit used
above.

Call a base prime missed when every source prime above its generic
point is $\pi$-exceptional.  The model property also makes all other
components of its pullback exceptional.  By Lemma~\ref{conv:return},
it is enough to construct a weighted metric satisfying
Theorem~\ref{jet:base-lemma} on $-K_Y+A$, with $A$ supported on missed
primes.  We now verify the only new ingredient, horizontal strictness,
using intersections on the fibers.

\subsection{Horizontal positivity from fiber intersections}

Let $\widetilde L=\pi^*L$ and let $\alpha$ be its smooth semipositive
curvature form.  Since $D$ is pseudoeffective and
$u\pi^*N_q-f^*A_0$ has an effective representative,
$k\widetilde L-f^*A_0$ is pseudoeffective for some $k>0$.
Let $J$ be ample on $S$, put $d_0=n-b$, and choose the largest
$t\in\{0,\ldots,d_0\}$ such that
\[
 (\widetilde L^tJ^{d_0-t}\cdot S_y)>0
\]
on a smooth general fiber.  The choice exists because $t=0$ works.
These intersections are constant on the smooth base locus.
Semipositivity then makes the restriction of $\alpha$ have rank at
most $t$ at every point of every such fiber.

The pseudoeffective comparison and nefness give
\[
 ((k\widetilde L)^b\cdot\widetilde L^tJ^{d_0-t})
 \ge ((f^*A_0)^b\cdot\widetilde L^tJ^{d_0-t})>0.
\]
Indeed the difference of the $b$-th powers factors as the
pseudoeffective class $k\widetilde L-f^*A_0$ times a sum of products
of nef classes.  The final strict inequality is the projection
formula and the definition of $t$.  The smooth representative shows
that $\alpha$ has rank at least $b+t$ on a nonempty open set.

A vertical null vector for a nonnegative Hermitian form is null for
the full form.  The vertical null space has dimension at least
$d_0-t$, so the total rank is at most $b+t$ on the smooth locus.
On the preceding open set equality holds, and the full kernel is
vertical.  A smaller constant-rank patch, disjoint from $E$, therefore
satisfies
\[
 \alpha\ge c f^*\omega_Y
\]
for some $c>0$ and a positive base form $\omega_Y$.

The adjoint rank is one by \eqref{jet:all-ranks}.
Proposition~\ref{coa:volume-metric} now gives an integral effective
$A$ supported on missed primes and a semipositive metric on
$-K_Y+A$, smooth and strictly positive on an ordinary open set,
with weighted density in $L^{1+\gamma}_{\rm loc}$ for some
$\gamma>0$.  Theorem~\ref{jet:base-lemma} makes $W+C_A$ rationally
effective with $C_A\ge0$ supported on $A$.
Lemma~\ref{conv:return} applied to $R$ gives a section of a positive
multiple of $L$.  This completes the jet-based proof of
Theorem~\ref{conv:fixed-error}.
